\documentclass{ximera}
\title{Cycloidal Arclengths: An Animated Approach, Abstract}

\newcommand{\pskip}{\vskip 0.1 in}

%\newtheorem{theorem19}{Robinson's Arclength Theorem for Prolate Cycloids, Local Version}
%\newtheorem{theorem20}{An Arclength Theorem for Prolate Epi and Hypocycloids}
%\newtheorem{theorem21}{Robinson's Lemma}
%\newtheorem{theorem22}{The Inclination Lemma}
%\newtheorem{theorem23}{An Arclength Theorem for General Cycloids}

\begin{document}
\begin{abstract}
Relationship between their arclengths.
\end{abstract}
\maketitle
 
\section{Introduction}

%Figure 1 shows a \emph{prolate cycloid} (from the Latin \emph{poferre}, to extend and the Greek \emph{kyklos}, circle). It is the curve traced by a point $P_1$ outside a circle rolling on a line

Roll a circle ${\cal C}_1$ of radius $a$ once through its circumference on a line ${\cal L}$ and a point $P_1$ fixed in the reference frame of the circle sweeps out one arch of a \emph{general cylcoid}. When $P_1$ is outside the circle, the \emph{prolate cylcoid} (from the Latin \emph{poferre}, to extend and the Greek \emph{kyklos}, circle) crosses ${\cal L}$. In a recent article,  D. Robinson showed the one-arch length of the prolate cycloid on the same side of ${\cal L}$ as ${\cal C}_1$ exceeds the opposite-side length by $8a$. The excess surprisingly does not depend on the distance $b>a$ of the tracing point from the center of ${\cal C}_1$.  Choosing $b=a$ shows the excess equals the one-arch length of the \emph{ordinary} $a$-cycloid traced by a point on the circumference of the same rolling circle (see Figure 1, where ${\cal L}$ is the $x$-axis, $a=2$, and $b=3.6$).

\begin{onlineOnly}
    \begin{center}
\desmos{ty0hdz7uyj}{900}{600}
\end{center}
\end{onlineOnly}

\href{https://www.desmos.com/calculator/ty0hdz7uyj}{Limacon Cylcoid Arclength 1B}

%What if $b<a$?. Then the tracing point lies inside the rolling circle and the \emph{curtate} (from the Latin \emph{curtus}, cut short) cycloid does not cross the $x$-axis. But it does have inflection points. And the convave-down arclength of one arch exceeds the concave-up length by $8b$ (Figure ??). Each of these arclength theorems follows quickly from the other. But it took us a long time to realize this and so we save this discussion for last (Section 7). %let alone to stumble upon the second theorem. So we save this discussion for last (Section 7).

If instead the rolling circle has radius $b$ and the tracing point lies $a<b$ units from its center,  the \emph{curtate} (from the Latin \emph{curtus}, cut short) cycloid does not cross the $x$-axis. But it does have inflection points. And the convave-down arclength of one arch exceeds the concave-up length by $8a$ (Figure ??). Each of these arclength theorems follows quickly from the other. But it took a long time for me to realize this, so we'll save this discussion for last (Section 7).

\begin{onlineOnly}
    \begin{center}
\desmos{0eosvsfm9r}{900}{600}
\end{center}
\end{onlineOnly}

\href{https://www.desmos.com/calculator/0eosvsfm9r}{Curtate Cycloid 33}

\end{document}
